\documentclass[12pt]{article}
\usepackage[utf8]{inputenc}
\usepackage[T1]{fontenc}
\usepackage{hyperref}
\title{Tai Pao Language Resources for K-12 Curriculum Development}
\author{Compiled by Claude Sonnet 5 (claude-sonnet-5) for Andr\'as Kornai}
\date{}
\begin{document}
\maketitle
\section{Introduction}
The code thc is retired. According to Glottolog, it was merged into tpo (Tai Pao) following a 2015 change request, adopted 2016-01-15; Tai Hang Tong is an alternate name for the same language. It is separate from Tai Hongjin (tiz). The Glottolog family label ``Bookkeeping'' in the task description reflects a placeholder classification for retired codes, not a real language family. No K-12 resources were found for this entry; any work should be done under tpo.
\section{Primary Educational Resources}
None found.
\section{Linguistic and Reference Materials}
\begin{itemize}
\item Glottolog record for ISO code thc: \url{https://glottolog.org/resource/languoid/iso/thc}
\item Wikipedia, Tai Pao language (current code tpo): \url{https://en.wikipedia.org/wiki/Tai_Pao_language}
\end{itemize}
\section{Religious and Community Resources}
None found.
\section{Digital and Audio Resources}
None found.
\section{Curriculum Development Context}
Because thc is a retired code, a dedicated curriculum cannot be based on it. The relevant language is Tai Pao (tpo), which should be researched separately.
\end{document}